\section{Conclusions}

This section consolidates the answers to the four research questions, identifies the policy conclusions that survive sensitivity analysis, and clarifies the external validation required before predictive use.

This study used an open, survey-informed ABM to evaluate how malicious news sources can imitate trusted-source attributes and how platform-controllable reach and recommendation policies interact with those strategies. Across 6,690 seeded runs in 423 conditions, source camouflage consistently improved audience-related outcomes, with combined imitation producing the largest gain. Campaign timing and duration determined how long the changed source could compete for attention. Immediate perfect truth signals strongly attenuated strategy effects, whereas incomplete, inaccurate or delayed recommendation evidence allowed substantially larger effects. Reach therefore cannot be interpreted independently of the recommendation mechanism that may reveal the source.

The sensitivity experiments provide the qualification needed for policy interpretation. Audience capture remained positive across all structural contrasts. False-exposure effects were less general: memory was the largest influence on combined imitation, several inputs interacted, and credibility camouflage alone seldom produced precise structural exposure effects. Consequently, policy evaluation should not equate malicious-source popularity with misinformation exposure, should use both source- and publication-level measures, and should test recommendation quality and reach governance jointly.

The model is a transparent computational laboratory, not a digital twin of X. Its numerical results are comparative outcomes under the stated model assumptions, not forecasts of X. Its value lies in controlled comparison, explicit assumptions and reproducible uncertainty analysis. External validation with observed X data is the next step required before predictive or population-level claims can be made.
